\chapter{Conclusion}
\label{ch:conclusion}

This thesis examined a practical question in interview-score modelling: does making a
predictor's scores more similar across learned visual or audio groups also reduce
disparity across independently supplied attributes, without losing useful prediction?
The answer requires separating predictive validity, proxy response and audited-group
transfer. None can be inferred from a small score-distribution distance alone.

The MIT experiment uses 138 recordings from 69 correctly paired participants. Its
outcome is crowd-perceived hireability, not employment success, and the low agreement
among individual raters remains a substantive measurement limit. The original neural
scorer is weaker than the ridge references. Within that design, adding a W1 penalty
to either matched sampling control lowers rank correlation, and none of the six
primary perceived-attribute disparity contrasts resolves a reduction. This is a
controlled negative result for the evaluated model, not a general rejection of proxy
regularisation.

The exploratory linear follow-up makes the mechanism test more informative. Its
task-only scorer achieves pooled SRCC 0.418 and normalised MAE 0.089 on the same
predictions for participants held out together across sessions. At penalty strength 0.02, the visual intervention
reduces mean held-out proxy W1 from 0.045 to 0.035; the corresponding audio intervention
reduces it from 0.087 to 0.068. Thus a working predictor can respond to the proxy
constraint. The full strength grid also exposes the predictive cost rather than
selecting a favourable point after inspecting the audit. All primary-strength fits
converge, while some strongest-penalty fits reach the iteration limit.

Transfer to the perceived attributes is less certain. The visual intervention lowers
the post-session perceived-gender W1 point estimate from 0.049 to 0.029, but its
paired reduction interval includes zero. Other appearance axes do not show a consistent
benefit. The annotations describe one observer's perception of particular recordings;
they do not establish demographic identity or inter-rater reliability. Visual proxy
partitions agree with this observer's gender-presentation codes in the held-out folds,
but that association is not a general demographic-classification result. The central
finding is therefore an active proxy mechanism with a conditional utility cost and
uncertain transfer across audited groups, not demonstrated hiring fairness.

AVI-Personality extends the evidence to independent participants and recruiter-rated
interview performance. A video/audio ridge predictor achieves test MAE 0.419 on its
1--5 scale, compared with 0.479 for the training-target mean. Its observed self-reported
male--female error difference is 0.067 rating points and remains descriptive.

A matched fifteen-arm comparison on those independent participants then supplies the
treatment-versus-control evidence the baseline alone could not. At the primary strength
the penalty changes original-scale error by 0.002, rank correlation by $-0.018$ and the
audited gender score gap by 0.001, with every interval containing zero. The independent
corpus therefore reproduces the main study's outcome under a different outcome construct
and rating source rather than revising it. Separating training from held-out proxy
response makes the reason legible: the penalty collapses the training visual proxy gap
almost entirely while leaving the held-out gap essentially unchanged, so the constraint is
satisfied on the participants used to fit it and does not describe those it is applied to.
The audio proxy constraint does transfer, reducing the held-out proxy gap by two thirds,
and still yields no demonstrable disparity benefit while incurring clear utility damage.
A proxy whose constraint generalises is therefore still not a usable stand-in for the
audited attribute. Because the comparison was specified after the external baseline had
been inspected, it is exploratory on independent people rather than confirmatory, and
different outcomes, full-frame rather than face-crop video and the absence of transcripts
continue to prevent an exact replication claim.

One limitation conditions all of these results and is stated here rather than left to be
inferred. Validation-selected shrinkage reached the upper limit of its search range in
every fold of every study reported here, leaving predicted scores with roughly half the
dispersion of the ratings they predict. A score-disparity penalty acts on the separation
between group prediction distributions, so a compressed predictor offers little separation
for it to remove. The absence of disparity movement is accordingly consistent with two
explanations this evidence cannot separate: that proxy regularisation does not reduce
audited disparity, or that the scorers available to it were too compressed for such a
reduction to be expressed. The learned configurations' failure to reach the deterministic
ridge floors, and the random-proxy control's equivalence to the structured proxies, both
favour the second as a contributing factor.

The contribution is a coherent empirical distinction between predicting a subjective
interview rating, manipulating a learned partition and improving a separately audited
group metric, together with evidence that the second does not deliver the third even when
it succeeds on held-out participants. Matched task-only controls, complete
penalty-strength reporting and the separation of training from held-out proxy response
make that distinction observable. Stronger evidence would require a scorer competitive
with its own deterministic floors, a regularisation search that does not terminate at its
boundary, more reliable measurement and replication across adequately supported groups.
The present findings justify those specific tests, while providing no basis for automated
hiring, causal coaching or employment-validity claims.
